\documentclass{article}
\usepackage{amsmath,amssymb}
\usepackage{algorithm}
\usepackage{algorithmic}

\begin{document}

\section{Monte Carlo Basics}


\section{Generalized Policy Iteration using Monte Carlo Methods}

\section{Monte Carlo Prediction}
\quad The objective of Monte Carlo prediction is to estimate the value function \( V_\pi(s) \) for a given policy \( \pi \) using sample episodes generated by following the policy. 
The key idea is to use the returns observed in these episodes to update the value estimates.


\subsection{On-Policy Monte Carlo Prediction}
\quad In on-policy Monte Carlo prediction, we generate episodes by following the policy \( \pi \) that we are trying to evaluate.
 However, when improving the policy, we can use an \(\epsilon\)-greedy approach to ensure that we explore the action space sufficiently. 
 This means that with probability \( 1 - \epsilon \), we follow the current policy, and with probability \( \epsilon \), we select a random action.

\subsection{Off-Policy Monte Carlo Prediction}
\quad In off-policy Monte Carlo prediction, we generate episodes using a different behavior policy \( b \) that may be more exploratory than the target policy \( \pi \) we want to evaluate. 
To correct for the discrepancy between the behavior policy and the target policy, we use importance sampling to weight the returns observed in the episodes. 

To obatin the importance sampling ratio, consider the following episode trajectory generated by the behavior policy \( b \):
\begin{equation}
S_0, A_0, S_1, A_1, \ldots, S_{T-1}, A_{T-1}, S_T
\end{equation}

This trajectory is a relaization of the behavior policy \( b \) and the probability of generating this trajectory under the behavior policy is given by:
\begin{equation}
P_b(S_0, A_0, S_1, A_1, \ldots, S_{T-1}, A_{T-1}, S_T) = P(S_0) \prod_{t=0}^{T-1} b(A_t | S_t) P(S_{t+1}| S_t, A_t)
\end{equation}

The probability of generating the same trajectory under the target policy \( \pi \) is given by:
\begin{equation}
P_\pi(S_0, A_0, S_1, A_1, \ldots, S_{T-1}, A_{T-1}, S_T) = P(S_0) \prod_{t=0}^{T-1} \pi(A_t | S_t) P(S_{t+1}| S_t, A_t)
\end{equation}

The importance sampling ratio \( \rho_t \) at time step \( t \) is defined as the ratio of the probability of generating the trajectory under the target policy to the probability of generating the trajectory under the behavior policy:
\begin{equation}
\rho_t = \prod_{k=0}^{t} \frac{\pi(A_k | S_k)}{b(A_k | S_k)}
\end{equation}  

Notice that the state transition probabilities \( P(S_{t+1}| S_t, A_t) \) and the initial state distribution \( P(S_0) \) cancel out in the ratio, which is a key advantage of using importance sampling in off-policy Monte Carlo prediction.

There are two common approaches to using importance sampling in off-policy Monte Carlo prediction: ordinary importance sampling and weighted importance sampling.

\textbf{Ordinary Importance Sampling:} In ordinary importance sampling, we use the importance sampling ratio \( \rho_t \) to weight the returns observed in the episodes and then simply average them.
The value estimate for state \( s \) is given by:
\begin{equation}
V(s) = \frac{1}{N(s)} \sum_{t=0}^{T-1} \rho_t G_t
\end{equation}
where \( N(s) \) is the number of times state \( s \) has been visited in the episodes.

\textbf{Weighted Importance Sampling:} In weighted importance sampling, we normalize the importance sampling ratios to ensure that they sum to 1. The value estimate for state \( s \) is updated as follows:
\begin{equation}
V(s) \leftarrow V(s) + \alpha \frac{\sum_{t=0}^{T-1} \rho_t (G_t - V(s))}{\sum_{t=0}^{T-1} \rho_t}
\end{equation}

\section{Monte Carlo Control}

\subsection{On-Policy Monte Carlo Control}

\subsection{Off-Policy Monte Carlo Control}
\quad In off-policy Monte Carlo control, we use a behavior policy \( b \) to generate episodes and a target policy \( \pi \) that we want to optimize.
The advantage of off-policy control is that we can use a more exploratory behavior policy to generate episodes, which can lead to better estimates of the value function and faster convergence to the optimal policy.



\end{document}